\section{Exercices d'application}

\rubrique{Ensemble de définition}

\exo{}\diff{1}
Déterminer l'ensemble de définition des fonctions suivantes :
\[
f(x)=\frac{2x+1}{x^2-4},\qquad g(x)=\sqrt{3x-6},\qquad h(x)=\frac{\sqrt{x+1}}{x-3}.
\]

\exo{}\diff{2}
Soit $f(x)=\frac{x^2-1}{x^2+2x-3}$. Déterminer $D_f$ puis simplifier l'expression de $f(x)$.

\exo{}\diff{2}
On considère $f(x)=\sqrt{4-x^2}+\frac{1}{\sqrt{x-1}}$. Déterminer $D_f$.

\rubrique{Parité et périodicité}

\exo{}\diff{1}
Étudier la parité des fonctions suivantes :
\[
f(x)=x^4-3x^2+1,\qquad g(x)=x^3-2x,\qquad h(x)=\frac{x}{x^2+1}.
\]

\exo{}\diff{2}
Soit $f$ définie sur $\R$ par $f(x)=x^2-2|x|+3$. Montrer que $f$ est paire. Tracer la courbe pour $x\ge 0$ sachant que $f(x)=x^2-2x+3$ sur $[0,+\infty[$.

\exo{}\diff{2}
Une fonction $f$ définie sur $\R$ vérifie $f(x+2)=f(x)$ pour tout $x\in\R$. On donne $f(x)=x^2$ pour $x\in[-1,1]$.
\begin{enumerate}[label=\alph*)]
\item Calculer $f(3)$ et $f(-4)$.
\item Tracer la courbe de $f$ sur $[-3,3]$.
\end{enumerate}

\rubrique{Sens de variation et taux de variation}

\exo{}\diff{1}
Soit $f(x)=3x-5$. Calculer le taux de variation de $f$ entre $1$ et $4$. Que peut-on en déduire ?

\exo{}\diff{2}
Soit $f(x)=x^2-4x+3$ sur $\R$.
\begin{enumerate}[label=\alph*)]
\item Calculer le taux de variation de $f$ entre $a$ et $b$.
\item En déduire le sens de variation de $f$ sur $]-\infty,2]$ et sur $[2,+\infty[$.
\end{enumerate}

\exo{}\diff{3}
Soit $f(x)=\frac{2x+1}{x-1}$ sur $\R\setminus\{1\}$.
\begin{enumerate}[label=\alph*)]
\item Calculer le taux de variation de $f$ entre $a$ et $b$ ($a\neq b$, $a\neq1$, $b\neq1$).
\item Montrer que $f$ est décroissante sur $]-\infty,1[$ et sur $]1,+\infty[$.
\end{enumerate}

\rubrique{Éléments de symétrie}

\exo{}\diff{2}
Soit $f(x)=x^2-4x+5$. Montrer que la droite $x=2$ est un axe de symétrie de la courbe de $f$.

\exo{}\diff{2}
Soit $f(x)=x^3-3x$. Montrer que le point $A(0,0)$ est un centre de symétrie de la courbe de $f$.

\exo{}\diff{3}
Soit $f(x)=\frac{x+1}{x-2}$. Déterminer le centre de symétrie de la courbe de $f$.

\rubrique{Fonctions associées et comparaison}

\exo{}\diff{1}
Soit $f(x)=x^2$. Tracer dans un même repère les courbes de $f$, $g(x)=f(x)+2$ et $h(x)=f(x+1)$.

\exo{}\diff{2}
Soit $f(x)=\sqrt{x}$ sur $[0,+\infty[$. Construire les courbes de $g(x)=f(x-1)$ et $h(x)=f(x)+1$.

\exo{}\diff{2}
Soit $f(x)=x^2-2x$. Déterminer le maximum de $f$ sur $[0,3]$.

\exo{}\diff{3}
Soit $f(x)=\frac{1}{x^2+1}$ sur $\R$. Montrer que $f$ est majorée par $1$ et minorée par $0$.

\section{Exercices d'entraînement}

\rubrique{Ensemble de définition}

\exo{}\diff{1} Déterminer $D_f$ pour $f(x)=\frac{3}{x-5}$.

\exo{}\diff{1} Déterminer $D_f$ pour $f(x)=\sqrt{2x+7}$.

\exo{}\diff{2} Déterminer $D_f$ pour $f(x)=\frac{\sqrt{x+2}}{x^2-1}$.

\exo{}\diff{2} Déterminer $D_f$ pour $f(x)=\sqrt{x^2-5x+6}$.

\exo{}\diff{3} Déterminer $D_f$ pour $f(x)=\frac{1}{\sqrt{9-x^2}}+\frac{1}{x-4}$.

\rubrique{Parité et périodicité}

\exo{}\diff{1} Étudier la parité de $f(x)=x^2+1$.

\exo{}\diff{1} Étudier la parité de $f(x)=x^3+x$.

\exo{}\diff{2} Soit $f(x)=|x|+x^4$. Montrer que $f$ est paire.

\exo{}\diff{2} Soit $f$ périodique de période $3$ et $f(x)=x$ sur $[0,3[$. Calculer $f(7)$ et $f(-2)$.

\exo{}\diff{3} Montrer que $f(x)=\cos x$ est paire et périodique de période $2\pi$.

\rubrique{Sens de variation}

\exo{}\diff{1} Étudier les variations de $f(x)=2x+1$ sur $\R$.

\exo{}\diff{2} Étudier les variations de $f(x)=x^2-2x$ sur $\R$.

\exo{}\diff{2} Étudier les variations de $f(x)=\frac{1}{x}$ sur $]0,+\infty[$.

\exo{}\diff{3} Soit $f(x)=\sqrt{x+1}$. Montrer que $f$ est croissante sur $[-1,+\infty[$.

\rubrique{Symétries}

\exo{}\diff{2} Montrer que $f(x)=x^2+2x+3$ admet un axe de symétrie.

\exo{}\diff{2} Montrer que $f(x)=x^3-3x+1$ admet un centre de symétrie.

\exo{}\diff{3} Déterminer le centre de symétrie de $f(x)=\frac{2x-1}{x+3}$.

\rubrique{Fonctions associées}

\exo{}\diff{1} Soit $f(x)=x^2$. Tracer $g(x)=f(x)-3$.

\exo{}\diff{2} Soit $f(x)=\sqrt{x}$. Tracer $g(x)=f(x+2)$.

\exo{}\diff{2} Soit $f(x)=|x|$. Tracer $g(x)=f(x-1)+2$.

\exo{}\diff{3} Soit $f(x)=x^2-4x$. Déterminer le minimum de $f$ sur $\R$.

\exo{}\diff{3} Montrer que $f(x)=\frac{2x}{x^2+1}$ est majorée par $1$ et minorée par $-1$.

\section{Problèmes type Probatoire/Bac}

\sujetbac[Probatoire Blanc — Étude d'une fonction rationnelle]
Soit $f$ la fonction définie par $f(x)=\frac{x^2-2x+3}{x-1}$.
\begin{enumerate}[label=\arabic*)]
\item Déterminer $D_f$.
\item Montrer que $f(x)=x-1+\frac{2}{x-1}$.
\item Étudier les variations de $f$ sur $]1,+\infty[$.
\item Montrer que le point $I(1,0)$ est un centre de symétrie de la courbe de $f$.
\item Tracer la courbe de $f$ sur $]1,+\infty[$.
\end{enumerate}

\sujetbac[Baccalauréat Blanc — Fonction avec valeur absolue]
Soit $f$ définie sur $\R$ par $f(x)=|x^2-4|$.
\begin{enumerate}[label=\arabic*)]
\item Exprimer $f$ sans valeur absolue.
\item Étudier la parité de $f$.
\item Étudier les variations de $f$ sur $[0,+\infty[$.
\item Tracer la courbe de $f$ sur $[-3,3]$.
\end{enumerate}

\sujetbac[Probatoire — Fonction et transformation]
Soit $f(x)=x^2-2x$ sur $\R$.
\begin{enumerate}[label=\arabic*)]
\item Déterminer les variations de $f$.
\item Soit $g(x)=f(x)+3$. Tracer $C_f$ et $C_g$.
\item Soit $h(x)=f(x-1)$. Tracer $C_h$.
\item Déterminer le minimum de $f$ sur $[-1,2]$.
\end{enumerate}

\section{Corrigés}

\corrige{de l'exercice 1}
\begin{enumerate}[label=\alph*)]
\item $f(x)=\frac{2x+1}{x^2-4}$. $D_f=\{x\in\R\mid x^2-4\neq0\}=\R\setminus\{-2,2\}$.
\item $g(x)=\sqrt{3x-6}$. $D_f=\{x\in\R\mid 3x-6\ge0\}=[2,+\infty[$.
\item $h(x)=\frac{\sqrt{x+1}}{x-3}$. Conditions : $x+1\ge0$ et $x-3\neq0$, soit $x\ge-1$ et $x\neq3$. Donc $D_f=[-1,3[\cup]3,+\infty[$.
\end{enumerate}

\corrige{de l'exercice 2}
$f(x)=\frac{x^2-1}{x^2+2x-3}$. Le dénominateur se factorise : $x^2+2x-3=(x+3)(x-1)$. Donc $D_f=\R\setminus\{-3,1\}$. Le numérateur : $x^2-1=(x-1)(x+1)$. Pour $x\neq1$, on simplifie : $f(x)=\frac{(x-1)(x+1)}{(x+1)(x-1)}=\frac{x+1}{x+3}$ (pour $x\neq-3$ et $x\neq1$).

\corrige{de l'exercice 3}
$f(x)=\sqrt{4-x^2}+\frac{1}{\sqrt{x-1}}$. Conditions : $4-x^2\ge0$ et $x-1>0$. $4-x^2\ge0\iff x^2\le4\iff -2\le x\le2$. $x-1>0\iff x>1$. Intersection : $]1,2]$. Donc $D_f=]1,2]$.

\corrige{de l'exercice 4}
\begin{enumerate}[label=\alph*)]
\item $f(x)=x^4-3x^2+1$. $f(-x)=(-x)^4-3(-x)^2+1=x^4-3x^2+1=f(x)$. Donc $f$ est paire.
\item $g(x)=x^3-2x$. $g(-x)=(-x)^3-2(-x)=-x^3+2x=-(x^3-2x)=-g(x)$. Donc $g$ est impaire.
\item $h(x)=\frac{x}{x^2+1}$. $h(-x)=\frac{-x}{(-x)^2+1}=-\frac{x}{x^2+1}=-h(x)$. Donc $h$ est impaire.
\end{enumerate}

\corrige{de l'exercice 5}
$f(x)=x^2-2|x|+3$. $f(-x)=(-x)^2-2|-x|+3=x^2-2|x|+3=f(x)$. Donc $f$ est paire. Pour $x\ge0$, $f(x)=x^2-2x+3=(x-1)^2+2$. La courbe pour $x\ge0$ est une parabole de sommet $(1,2)$.

\corrige{de l'exercice 6}
$f$ est périodique de période $2$.
\begin{enumerate}[label=\alph*)]
\item $f(3)=f(3-2)=f(1)=1^2=1$. $f(-4)=f(-4+2\times2)=f(0)=0$.
\item Sur $[-1,1]$, $f(x)=x^2$. Par périodicité, on reproduit ce motif sur $[-3,-1]$, $[1,3]$, etc.
\end{enumerate}

\corrige{de l'exercice 7}
$f(x)=3x-5$. Taux de variation entre $1$ et $4$ : $\frac{f(4)-f(1)}{4-1}=\frac{(12-5)-(3-5)}{3}=\frac{7-(-2)}{3}=\frac{9}{3}=3>0$. Donc $f$ est croissante sur $\R$.

\corrige{de l'exercice 8}
\begin{enumerate}[label=\alph*)]
\item $f(x)=x^2-4x+3$. Taux entre $a$ et $b$ : $\frac{f(b)-f(a)}{b-a}=\frac{(b^2-4b+3)-(a^2-4a+3)}{b-a}=\frac{b^2-a^2-4(b-a)}{b-a}=a+b-4$.
\item Pour $a,b\in]-\infty,2]$, $a+b\le4$, donc $a+b-4\le0$ : $f$ est décroissante. Pour $a,b\in[2,+\infty[$, $a+b\ge4$, donc $a+b-4\ge0$ : $f$ est croissante.
\end{enumerate}

\corrige{de l'exercice 9}
\begin{enumerate}[label=\alph*)]
\item $f(x)=\frac{2x+1}{x-1}$. Taux : $\frac{f(b)-f(a)}{b-a}=\frac{\frac{2b+1}{b-1}-\frac{2a+1}{a-1}}{b-a}=\frac{(2b+1)(a-1)-(2a+1)(b-1)}{(b-a)(b-1)(a-1)}=\frac{-3}{(a-1)(b-1)}$.
\item Sur $]-\infty,1[$, $a-1<0$, $b-1<0$, donc $(a-1)(b-1)>0$, le taux est négatif : $f$ décroissante. Sur $]1,+\infty[$, $a-1>0$, $b-1>0$, même conclusion.
\end{enumerate}

\corrige{de l'exercice 10}
$f(x)=x^2-4x+5=(x-2)^2+1$. Pour tout $h$, $f(2+h)=(2+h)^2-4(2+h)+5=4+4h+h^2-8-4h+5=h^2+1$. $f(2-h)=(2-h)^2-4(2-h)+5=4-4h+h^2-8+4h+5=h^2+1$. Donc $f(2+h)=f(2-h)$ : la droite $x=2$ est axe de symétrie.

\corrige{de l'exercice 11}
$f(x)=x^3-3x$. $f(-x)=(-x)^3-3(-x)=-x^3+3x=-(x^3-3x)=-f(x)$. $f$ est impaire, donc sa courbe est symétrique par rapport à l'origine $O(0,0)$.

\corrige{de l'exercice 12}
$f(x)=\frac{x+1}{x-2}=1+\frac{3}{x-2}$. On pose $X=x-2$, $Y=y-1$. Alors $Y=\frac{3}{X}$. La fonction $X\mapsto\frac{3}{X}$ est impaire, donc le point $I(2,1)$ est centre de symétrie.

\corrige{de l'exercice 13}
$f(x)=x^2$ (parabole sommet $O$). $g(x)=f(x)+2=x^2+2$ (translation verticale de $+2$). $h(x)=f(x+1)=(x+1)^2$ (translation horizontale de $-1$).

\corrige{de l'exercice 14}
$f(x)=\sqrt{x}$ (demi-parabole). $g(x)=f(x-1)=\sqrt{x-1}$ (translation de $+1$ en $x$). $h(x)=f(x)+1=\sqrt{x}+1$ (translation de $+1$ en $y$).

\corrige{de l'exercice 15}
$f(x)=x^2-2x=(x-1)^2-1$. Sur $[0,3]$, le sommet est en $x=1$ avec $f(1)=-1$. $f(0)=0$, $f(3)=3$. Le maximum est $3$ atteint en $x=3$.

\corrige{de l'exercice 16}
$f(x)=\frac{1}{x^2+1}$. $x^2+1\ge1$, donc $0<\frac{1}{x^2+1}\le1$. Donc $f$ est minorée par $0$ (non atteint) et majorée par $1$ (atteint en $x=0$).

\corrige{du problème 1}
\begin{enumerate}[label=\arabic*)]
\item $f(x)=\frac{x^2-2x+3}{x-1}$. $D_f=\R\setminus\{1\}$.
\item $x-1+\frac{2}{x-1}=\frac{(x-1)^2+2}{x-1}=\frac{x^2-2x+1+2}{x-1}=f(x)$.
\item Sur $]1,+\infty[$, $x-1>0$. $f(x)=x-1+\frac{2}{x-1}$. Taux de variation : $\frac{f(b)-f(a)}{b-a}=1-\frac{2}{(a-1)(b-1)}$. Pour $a,b>1$, $(a-1)(b-1)>0$, donc le taux peut être positif ou négatif. Étudions la dérivée (hors programme) ou utilisons la forme : $f$ est décroissante sur $]1,2]$ et croissante sur $[2,+\infty[$ (car $x-1$ croît et $\frac{2}{x-1}$ décroît).
\item $f(1+h)=h+\frac{2}{h}$, $f(1-h)=-h-\frac{2}{h}=-(h+\frac{2}{h})=-f(1+h)$. Donc $I(1,0)$ est centre de symétrie.
\item Tracé : asymptote verticale $x=1$, asymptote oblique $y=x-1$, minimum en $x=2$ ($f(2)=3$).
\end{enumerate}

\corrige{du problème 2}
\begin{enumerate}[label=\arabic*)]
\item $f(x)=|x^2-4|=\begin{cases} x^2-4 & \text{si } x\le-2 \text{ ou } x\ge2 \\ 4-x^2 & \text{si } -2\le x\le2 \end{cases}$.
\item $f(-x)=|(-x)^2-4|=|x^2-4|=f(x)$. Donc $f$ est paire.
\item Sur $[0,+\infty[$, $f(x)=\begin{cases} 4-x^2 & \text{sur } [0,2] \\ x^2-4 & \text{sur } [2,+\infty[ \end{cases}$. $f$ est décroissante sur $[0,2]$ et croissante sur $[2,+\infty[$.
\item Tracé : paraboles raccordées en $x=2$.
\end{enumerate}

\corrige{du problème 3}
\begin{enumerate}[label=\arabic*)]
\item $f(x)=x^2-2x=(x-1)^2-1$. Décroissante sur $]-\infty,1]$, croissante sur $[1,+\infty[$.
\item $g(x)=f(x)+3=x^2-2x+3=(x-1)^2+2$. Translation verticale de $+3$.
\item $h(x)=f(x-1)=(x-1)^2-2(x-1)=x^2-4x+3=(x-2)^2-1$. Translation horizontale de $+1$.
\item Sur $[-1,2]$, $f(-1)=3$, $f(1)=-1$, $f(2)=0$. Le minimum est $-1$ en $x=1$.
\end{enumerate}